\section{Problem formulation}
\label{sec:problem}

\subsection{A validation test as a decision rule}
\label{sec:pf-rule}

A validation test $T$ takes the data $Y$ of a validation sample, such as the default counts of a portfolio, the scores and outcomes of its obligors, or two samples of a model input, and returns a p-value $p_T(Y)\in[0,1]$, or no p-value when its statistic is undefined, for example when a variance is zero.
At a level $\lev$ the test rejects when the p-value falls below it, $p_T(Y)<\lev$; the engine that is measured colours a result amber below 5\,\% and red below 1\,\%, and the two levels of the measurement are these, $\lev\in\{0.05,0.01\}$.
A result without a p-value is never a rejection and is counted separately.
Let $G_\theta$ be a generator of validation samples with parameter $\theta$ and exactly known truth, and let the \emph{rejection probability} of the test at level $\lev$ be
\begin{equation}
\rej_T(\theta,\lev)\;=\;\Prob_\theta\big(p_T(Y)<\lev\big),\qquad Y\sim G_\theta .
\label{eq:rejection}
\end{equation}
Every defect of a model is a departure of $\theta$ from the set $\Theta_0(T)$ of the parameters for which the null hypothesis of $T$ holds. Some tests are one-sided, as the binomial test whose alternative is a probability of default above the one applied, and some are two-sided, as the DeLong comparison of two areas; the rejection rule is the same for both, and only the null set differs.

\subsection{Size, power and the estimands}
\label{sec:pf-estimands}

The size of the test is the largest rejection probability over the null set, and its power at a parameter outside it is the rejection probability there:
\begin{equation}
\mathrm{size}_T(\lev)\;=\;\sup_{\theta\in\Theta_0(T)}\rej_T(\theta,\lev),\qquad
\mathrm{power}_T(\theta,\lev)\;=\;\rej_T(\theta,\lev),\quad \theta\notin\Theta_0(T).
\label{eq:size-power}
\end{equation}
For the three tests whose rejection region is a single default count, the null is one-sided in the probability of default, along which the rejection probability does not decrease (\cref{prop:monotone}); at the correlation that the test assumes, the supremum is then attained at the boundary of the null, the generator $G_{\theta_b}$ in which the null holds with equality and every assumption of the test is met, and the size is $\rej_T(\theta_b,\lev)$.
For the other eighteen tests the null set contains nuisance parameters: any common distribution of two samples (the stability indexes and the tests of homogeneity), any vector of probabilities of default (the chi-square, Hosmer--Lemeshow and Spiegelhalter tests), any path of yearly probabilities (the normal and traffic-lights tests), any correlation between two scores (DeLong), any common value of the area, the concentration or the migration frequency that is compared (the discrimination, concentration and migration tests), any transition chain (the four tests of transitions).
For these the measurement does not take the supremum: it measures the rejection probability $\rej_T(\theta_b,\lev)$ at one null configuration, which is a lower bound on the size, and the statement that a test holds its size at that configuration says nothing of the other configurations of the null (the Jeffreys case in \cref{sec:li-external} shows how much a rate can move with the configuration).
\emph{Evidence} about $\rej_T(\theta_b,\lev)$ is its exact value where the test has one, or a rate $\hat\rej$ in $\nrep$ repetitions with its Wilson interval $[\ell,u]$ of \eqref{eq:se-wilson}, and the evidence \emph{shows an excess} at the level $\lev$ when the exact value is above $\lev$ or when $\ell>\lev$.
\begin{definition}[Holding a size]
\label{def:holds}
A test \emph{holds its size on a body of evidence} when the evidence shows no excess at either level, $\lev\in\{0.05,0.01\}$. It is \emph{conservative} at a level when the exact value, or the whole interval, lies below the level, and it \emph{does not hold its size} when the evidence shows an excess at one of the levels.
\end{definition}
Every count of tests that hold their size in this paper, in the title, the abstract, the results and the conclusions, applies \cref{def:holds} and names its evidence. The count is a statement about the evidence and not only about the test: a test that holds its size on a cell of 4,000 repetitions may exceed its level by less than the resolution of the cell (\cref{prop:resolution}), and a test shown above its level by 200,000 repetitions may exceed it by a small fraction of a point. The one-sided bound of \cref{sec:pm-judgement}, which the pipeline applies to every cell and which tolerates a larger excess, is reported beside the definition under the name of the screen, never instead of it.

A defect is a one-parameter family $\theta(\sev)$ of generators, with a severity $\sev$ on a ladder of six values, and the power curve of the test against it is $\sev\mapsto\rej_T(\theta(\sev),\lev)$.
The summary of a curve that the measurement reports is the severity at which the power first reaches 80\,\%,
\begin{equation}
\sev_{0.8}(T)\;=\;\inf\{\sev:\ \rej_T(\theta(\sev),0.05)\ge 0.8\},
\label{eq:s80}
\end{equation}
which is a property of the test and the ladder together and carries the unit of the ladder.
The estimands are therefore the size at 5\,\% and at 1\,\% of each test at the boundary of its null, its power at each point of the ladder of each defect that it targets, and $\sev_{0.8}$; for the three tests whose rejection region is a single default count, they are also computable exactly.

\subsection{The tests that are measured}
\label{sec:pf-tests}

The battery is the 21 tests of the five credit families of the engine that return a p-value: nine tests of the calibration of the default probability, two of discrimination, four of population and characteristic stability, two of the rating system, and four of rating transitions.
\Cref{tab:tests} lists each with its statistic and its null hypothesis, as the engine documents them.
The four transition tests have a size in the measurement and no power, because no ladder of a defect of the transition process is planted for them, and the measurement does not claim a power for them. The 21 tests of the other seven families of the engine are measured separately (\cref{sec:extension}), each at the null of a named study and at one stated alternative.

\begin{table}[htbp]
\centering
\footnotesize
\caption{The 21 tests of the credit families that are measured: the statistic each computes and the null hypothesis it tests, as documented by the engine; D is the number of defaults, N the number of obligors, and a grade is a rating class.}
\label{tab:tests}
\fitwidth{\begin{tabular}{p{3.0cm}p{5.0cm}p{6.4cm}}
\toprule
Test & Statistic and reference distribution & Null hypothesis \\
\midrule
Binomial & upper tail of $\Bin(N,\mathrm{PD})$ at $D$ & true PD not above the PD applied, defaults independent \\
Jeffreys & posterior $\mathrm{Beta}(D+\tfrac12,N-D+\tfrac12)$ at the PD applied & PD applied not below the true PD \\
Vasicek-corrected binomial & binomial mixture over one systematic factor, assumed asset correlation & true PD not above the PD applied at the assumed correlation \\
Chi-square over grades & $\sum_i (n_ip_i-d_i)^2/(n_ip_i(1-p_i))$, $\chi^2$ with one degree of freedom per grade & every grade PD is the true one, defaults independent \\
Default profile & Pearson statistic of the defaults over grades given their number, Monte Carlo p-value & defaults fall on the grades as the PD curve implies (the shape, not the level) \\
Hosmer--Lemeshow & groups by predicted PD, $\chi^2$ with $G$ degrees of freedom for given PDs & predicted equals true PD in every group \\
Spiegelhalter & standardised Brier score, normal tail & the forecast equals the true PD of every obligor \\
Normal, five years & $\sum_t(d_t-\mathrm{PD}_t)/(\sqrt T\hat\sigma)$, normal tail & no year's true PD exceeds its forecast, years independent \\
Traffic lights & multinomial profile of the colours of the yearly standardised default rates, exact enumeration & as the normal test, with the colour probabilities of the rule \\
AUC against development & $(\mathrm{AUC}_\mathrm{dev}-\mathrm{AUC}_\mathrm{cur})/s$, normal tail, development AUC taken as known & AUC at development smaller than in the observation period \\
DeLong & difference of two AUCs on one sample over its covariance-based standard error, two-sided & the two AUCs are equal \\
Population stability index & PSI against $(1/n+1/m)\chi^2_{1-\lev,B-1}$ (chi-square benchmark) & the binned distributions of the two samples are equal \\
Characteristic stability index & the same on one model input & the binned distributions of the input are equal \\
Chi-square of homogeneity & Pearson statistic of the $2\times B$ table & the binned distributions of the two samples are equal \\
Kolmogorov--Smirnov & supremum distance of the two empirical distributions & both samples come from one continuous distribution \\
Migration z-tests & one z per cell against its closer neighbour, smallest Holm-adjusted p & migration frequencies do not increase away from the diagonal \\
Concentration (ECB) & $\sqrt{K-1}(\mathrm{CV}_\mathrm{cur}-\mathrm{CV}_\mathrm{dev})/\sqrt{\mathrm{CV}_\mathrm{cur}^2(\tfrac12+\mathrm{CV}_\mathrm{cur}^2)}$, normal tail & concentration in the grades has not increased since development \\
Time homogeneity & chi-square on the transition counts of the periods & the transition probabilities are the same in every period \\
Markov order & chi-square on triplets of ratings grouped by direction & the rating process is of first order \\
Matrix against a reference & chi-square of the counts of a year against a reference matrix & the transition matrix equals the reference \\
Downward momentum & hazard-based test of a momentum coefficient & the momentum coefficient is zero \\
\bottomrule
\end{tabular}}
\end{table}

\subsection{Notation and assumptions}
\label{sec:pf-notation}

\Cref{tab:notation} collects the symbols.
The assumptions of the measurement are three.
The repetitions of one cell are independent draws of the generator, which the stream design secures.
The generator reproduces the null or the defect it states, which the exact probabilities check for the count tests and the reruns of published studies check for others.
The severity ladders are design choices, set so that the power of the targeted tests runs from their size towards one in the sample size of the variant; a power measured on a ladder is a property of the pair of test and ladder, and the limitations return to what that does not mean. The notation separates the probability $\rej$, a property of a test and a generator, from its estimate $\hat\rej$, a property of a sample of repetitions, because the conclusions of the text are statements about the first that are read from the second.

\begin{table}[htbp]
\centering
\small
\caption{Notation of the measurement: each symbol with its meaning; every quantity of the text is a probability or a rate unless a unit is stated, and a level is written as a fraction or as a percentage.}
\label{tab:notation}
\fitwidth{\begin{tabular}{p{2.6cm}p{11.6cm}}
\toprule
Symbol & Meaning \\
\midrule
$T$, $p_T(Y)$ & a validation test and its p-value on data $Y$ \\
$\lev$ & level of the test, 0.05 or 0.01 \\
$G_\theta$, $\theta_b$ & generator of known truth with parameter $\theta$; the parameter at the boundary of the null \\
$\rej_T(\theta,\lev)$ & rejection probability of $T$ at level $\lev$ under $G_\theta$ \\
$\nrep$, $K$, $\hat\rej=K/\nrep$ & number of repetitions of a cell, number of rejections, the measured rate \\
$\MCSE(\hat\rej)$ & Monte Carlo standard error of the rate, $\sqrt{\hat\rej(1-\hat\rej)/\nrep}$ \\
$z_q$ & the $q$ quantile of the standard normal distribution \\
$b_\nrep(\lev)$ & one-sided bound of the size judgement at $\nrep$ repetitions \\
$k^*$ & smallest default count whose p-value is below $\lev$ \\
$\sev$, $\sev_{0.8}$ & severity on a ladder; severity at which the power reaches 80\,\% \\
$\rho$, $\mathrm{PD}$ & asset correlation of the one-factor model; probability of default \\
$\delta_\nrep(\lev)$ & excess over the level that the bound detects with probability 0.8 \\
\bottomrule
\end{tabular}}
\end{table}
